\section{What result are we seeking?}

Installment 6: “Folkebladet,” March 29, 1916

And here no one among us should shrink back from the labor and exertion that such preparation requires. For it repays us well in our own spiritual development and in our fitness for battle against the enemy. If we truly desire that our participation in the Lord’s fighting army should yield some result, should be something other and more than a “beating of the air,” then preparation is necessary—regardless of what it costs in toil, labor, and prayer. I know very well that we can obtain recognition, and perhaps even praise, from some people without laboring with great exertion to find our own way into the word and without seeking to lead our hearers into it. For what some seek to obtain as the yield of a meeting is that the emotional life of the soul should in some way be set in motion. And the preacher who speaks in such a manner that this is attained often receives praise from superficial persons. But everyone knows that such results can easily be attained without preparation and without prayer—and many are tempted to be satisfied with such things. This, however, is not, and should never be, what we aim at when we participate in a meeting by bearing testimony. Human praise or human blame should have little influence upon us. Our calling is to set our honor in “holding Jesus Christ forth to view,” whether this brings us censure or recognition. For what is worth more than human praise or human blame is to receive a friendly, approving glance from him, our great and exalted king, in whose army we have enlisted and under whose banner of the cross we fight.

But his approval is never gained by fighting with self-made weapons, nor by seeking honor for ourselves.

Another matter to which I would also gladly point, because in my judgment it is so significant, is that the person who participates in a meeting should have an aim and a plan for his testimony: an aim, so that he intends something, and a plan, so that he pursues his aim.

We need not listen long to the testimony of some before it becomes clear to us that, at bottom, they aim at nothing and particularly desire nothing, except perhaps this: to speak for the sake of speaking.

That this is wrong should be evident to everyone. And the person who intends nothing by his speech should never occupy the meeting’s time. There are others who speak, and in whose testimony it can plainly be perceived that they do intend something. There are spirit and fire in them.

But what do they intend?

Some of them are almost offended by such a question and answer with indignation: What do we intend? We intend awakening. We intend that the people to whom we speak should be awakened. Yes, is this, then, not right? Certainly it is. And I rejoice every time that, in a meeting, I am permitted to hear a testimony borne forth by a heart filled with a burning desire that those who sleep may be awakened.

For awakening is continually needed, so long as the evil powers in the world wind their creeping courses around human beings.

The individual and the Lord’s congregation need awakening—such seasons of breakthrough in which life is renewed and made effective. And I will never join those who would dampen the zeal of those whose testimony aims at awakening.

No, let the word of awakening sound forth with clarity and power in our meetings! For the winter cold and the sleepiness are easily perceived everywhere. And when those who long for the manifestations of spring’s life, for the disappearance of the night—those whose eyes are weary from watching for the morning and wet with tears from sorrow over themselves and their brothers and sisters in the world—when such persons, with burning and throbbing hearts, speak for awakening in a meeting, then I rejoice. And from my innermost heart I desire that God may bless all those who are grieved over the spiritual distress among our people, and who therefore lament and long, pray and labor, testify and act, that awakening may come! For as a rule it comes only through labor, sacrifice, and prayer. But even if it comes, even if—as I have already pointed out—there arises some more or less powerful spiritual movement, and even if one or more persons indicate that they desire intercessory prayer and guidance, it is not enough. And the testimony that aims only at awakening, only at movement, is imperfect and one-sided.

Spiritual movement, where it is genuine, is only a promising beginning, and nothing more. The apple blossom in spring is beautiful; that is true. But the blossom is not thereby a guarantee of fruit. What is decisive is the summer and the harvest. Then it will become evident whether the blossom became more than a promising possibility. So it is also in the world of the Spirit. Awakening may be promising, but it is not the same as conversion and life in God.

But at times the spiritual stirrings and movements produced by a testimony in a meeting are not even genuine. The movement may be only a superficial intoxication that results in nothing, or in something worse than nothing: it makes the person harder and less receptive to the truth than before. It is therefore so important that the person who speaks for awakening in a meeting should set as his aim that the souls whom he seeks to reach through his testimony should be affected not only in their emotional life, but also in their volitional life, so that, through the working of God’s Spirit by means of the word, there may be conversion, change of mind, and new life in God. Only then has the testimony produced the effect that it ought to have.

It has been shown, not once but many times, that great movements in a meeting fulfilled very little of what they promised. The change of mind that appeared to be traceable in several persons lasted, in many cases, no longer than until the end of the emotionally charged meeting. And then it is certain that the very testimony that was spoken was not what it ought to have been.

The Lord’s word now, as before, has power in itself to recreate human hearts, so that it may truthfully be said of them: “The old has passed away; behold, all has become new.”

And where the Lord’s truth is spoken by living, Spirit-filled witnesses in a meeting prepared and conducted under earnest prayer, it will always bear its blessed fruit, even though no greater awakening is perceived.

[Page 291]

For the highest fruit a meeting can bear is that believers are awakened. Someone may perhaps say that believers are awake. Yes, perhaps they are. But they can fall asleep again: “When the bridegroom delayed, they all slumbered and fell asleep.” And therefore they continually need to be awakened further—awakened to see their own dullness and the great spiritual peril of their fellow human beings. It may indeed be that the greatest obstacle to the kingdom of God, to “the salvation of souls,” is the half-hearted, powerless godliness of those who confess Christ. Is it not a truth established by experience, yet deeply humiliating, that our dear Savior has often offered us an opportunity to “take up the cross,” but—we declined?

It did not escape the attention of him whose solicitous gaze is fixed upon us and who asks with concern: “Will you also go away?”

And perhaps it did not escape the attention of those who expect us to be “the light of the world,” and who can be won for Jesus Christ only through the testimony given by the life of God’s congregation—the life lived in voluntary renunciation and suffering, and in continual, self-forgetting labor for his cause. Christ has never, through the baptism of blood upon the cross, given his life for the salvation of the world so that we should discuss religious questions, attend festivals or meetings in order to enjoy ourselves, and otherwise remain as inconspicuous as possible—“laying up treasures upon earth.”

No, what is needed in order that the fire may be kindled of which Jesus said, “Fire I have come to cast upon the earth, and how gladly I would have it already kindled,” is that God’s congregation should awaken—should see its own danger and the danger of others, and should be penetrated throughout by an “unceasing pain” over the distress of our people and by a glad, battle-ready spirit of conquest whose aim is to win our people for Christ. And where such results are attained through a meeting, there God’s Spirit, who is the true power able to produce awakening, has been at work through the word, even though no tumultuous movement could be perceived.

For the awakening of believers most commonly takes place in such a manner—at the beginning, at least—that no one other than the Lord and the individual soul awakened by the word’s sharp rebuke knows anything of it.

This will, of course, soon be perceived by others as well, since the believer’s outward life and work will bear the mark of the renewal and enrichment of godliness.